% !TeX program = xelatex
% UTF-8 transcription; source wording and numbered/arrow-linked notes retained.
% Standalone XeLaTeX document; may also be input after loading xepersian.
% Original images (unmodified), in handwritten page order:
% images/pic01.jpg = 20261016_223739123.jpg
% images/pic02.jpg = 20261016_223746881.jpg
% images/pic03.jpg = 20261016_223752176.jpg
% images/pic04.jpg = 20261016_223757941.jpg
% Printed stationery text is retained in the page headers and footers.
\ifdefined\settextfont
  \def\RefractiveStandalone{0}
\else
  \def\RefractiveStandalone{1}
  \documentclass[12pt,a4paper]{article}
  \usepackage[margin=2cm]{geometry}
  \usepackage{amsmath,amssymb}
  \usepackage{xepersian}
  \settextfont{Amiri}
  \setlatintextfont{FreeSerif}
\fi
\newcommand{\refen}[1]{\lr{#1}}
\newcommand{\refunclear}[1]{#1\footnote{خوانش نامطمئن دست‌خط؛ تصویر اصلی بررسی شود.}}
\newcommand{\refpage}[1]{%
  \clearpage
  \noindent\refen{Note book}\hfill\refen{Subject:} ۱۱\par
  \noindent\refen{Veor:}\hfill\refen{Month:}\hfill\refen{Date:}\par
  \smallskip\hrule\medskip
  \noindent\textbf{صفحهٔ #1}\par\medskip}
\newcommand{\reffooter}[1]{\par\vfill\noindent
  \refen{BAHARE DANESH ASL}\hfill #1\par}
\ifnum\RefractiveStandalone=1 \begin{document}\fi

\refpage{۱}
\noindent\refen{Opthalmology}
\section*{عیوب انکساری چشم}

\subsection*{۱ ـ قرنیه و لنز ساختارهای انکساری چشم هستند.}
\begin{itemize}
  \item قرنیه \refen{$2/3$} و لنز \refen{$1/3$} قدرت انکساری چشم.
\end{itemize}

\subsection*{۲ ـ کاهش حدت بینایی}
\begin{itemize}
  \item[\refen{a)}] هایپروپی یا دوربینی:
    طول آگزیال کرهٔ چشم خیلی کوتاه باشد.
  \item[\refen{b)}] میوپی یا نزدیک‌بینی:
    طول آگزیال کرهٔ چشم خیلی بلند باشد.
  \item[\refen{c)}] آستیگماتیسم:
    قدرت انکساری قرنیه و لنز در یک محور نسبت به محور دیگر متفاوت باشد.
\end{itemize}

\subsection*{۳ ـ \refen{Pinhole}}
\begin{itemize}
  \item \refunclear{قطر مؤثر} مردمک خیلی باریک می‌شود.
  \item تاری دید ناشی از عیب انکساری به حداقل می‌رسد.
  \item با کمک آن می‌توان تخمین زد که حدت بینایی پس از اصلاح عیب انکساری چقدر خواهد بود.
\end{itemize}
\reffooter{۱}

\refpage{۲}
\subsection*{۴ ـ علل میوپی یا نزدیک‌بینی}
\begin{itemize}
  \item بزرگ‌تر بودن غیرطبیعی طول محور قدامی خلفی کرهٔ چشم:
    میوپی محوری.
  \item افزایش قدرت اجزای انکساری چشم:
    میوپی رفراکتیو.
  \item تغییر در اندکس انکساری اجزای داخل چشم:
    میوپی ضریبی.
\end{itemize}

\subsection*{۵ ـ امتروپیا ـ \refen{Emmetropia}}
وضعیتی که در آن هیچ‌گونه عیب انکساری در چشم وجود ندارد و اشعهٔ نور موازی با محور بینایی بدون تطابق روی فووآ متمرکز می‌گردد.

\subsection*{۶ ـ آمتروپیا}
عیب انکساری وجود دارد.
\begin{itemize}
  \item انواع: میوپی، هایپروپی، آستیگماتیسم.
\end{itemize}

\subsection*{۷ ـ آنیزومتروپی}
اختلاف عیب انکساری بین دو چشم یک فرد.

\subsection*{۸ ـ اصلاح میوپی}
استفاده از عدسی مقعر.
% A separate handwritten numeral 6 appears below the printed footer on pic02.
\par\vfill\noindent ۶
\reffooter{۲}

\refpage{۳}
\subsection*{۹ ـ نزدیک‌بینی محوری (آگزیال)}
\begin{itemize}
  \item شایع‌ترین فرم نزدیک‌بینی.
  \item انحنای عدسی و قرنیه طبیعی است.
  \item به ازای \refen{1 mm} افزایش طول محور چشم، ۳ دیوپتر نزدیک‌بینی ایجاد می‌شود.
\end{itemize}

\subsection*{۱۰ ـ نزدیک‌بینی انکساری (رفراکتیو)}
\begin{itemize}
  \item مهم‌ترین عامل، شعاع انحنای قرنیه است.
  \item طبیعی: \refen{7.8 mm}.
  \item افزایش انحنای قرنیه (قوز قرنیه): \refen{Keratoconus}.
\end{itemize}

\subsection*{۱۱ ـ میوپی ضریبی (\refen{Index})}
\begin{itemize}
  \item مراحل اولیهٔ آب مروارید.
  \item \refen{Second Sight}.
\end{itemize}

\subsection*{۱۲ ـ عوارض میوپی}
\begin{itemize}
  \item تغییرات ته چشم.
  \item آتروفی منتشر کوروئید و رتین در اطراف دیسک: مهم‌ترین.
  \item استافیلوما‌ی خلفی.
  \item مناطق آتروفیک در ناحیهٔ ماکولا.
  \item افزایش احتمال جداشدگی شبکیه: خطرناک.
\end{itemize}
\reffooter{۳}

\refpage{۴}
\begin{itemize}
  \item[$\ast$] آتروفی کوروئید در میوپی پاتولوژیک دیده می‌شود.
    \begin{itemize}
      \item بالاتر از \refen{6 D}.
    \end{itemize}
  \item تغییرات در میدان بینایی:
    \begin{itemize}
      \item در نزدیک‌بین‌ها شدید.
      \item \refen{DDX}: گلوکوم.
    \end{itemize}
  \item[$\ast$] بیماران با میوپی پاتولوژیک از نظر بروز گلوکوم و تغییرات شبکیه باید پیگیری شوند.
\end{itemize}
\reffooter{۴}

\ifnum\RefractiveStandalone=1
  \let\refend\enddocument
\else
  \let\refend\relax
\fi
\refend
